\chapter{Il tessuto epiteliale}

% ---------------------------------------------------------------------------
\section{I tessuti del corpo umano}

I \textbf{tessuti} sono gruppi di cellule che lavorano in cooperazione e che hanno simili caratteristiche morfologiche. Permettono una ``divisione del lavoro'' fra le cellule dell'organismo. Nel corpo umano si riconoscono \textbf{4 tipi di tessuti}:

\begin{itemize}
  \item \textbf{Epiteliale}: copre superfici esposte, delimita dotti e cavità interne, produce secreti ghiandolari.
  \item \textbf{Connettivo}: riempie gli spazi interni, offre supporto strutturale, conserva energia.
  \item \textbf{Muscolare}: si contrae per produrre un movimento attivo.
  \item \textbf{Nervoso}: conduce impulsi elettrici e trasporta informazioni.
\end{itemize}

L'organizzazione gerarchica del corpo umano procede dalle molecole alle cellule, dai tessuti agli organi, fino agli apparati. I tessuti con funzioni speciali si uniscono a formare organi, che interagiscono negli apparati.

% ---------------------------------------------------------------------------
\section{Caratteristiche generali del tessuto epiteliale}

Il \textbf{tessuto epiteliale} è formato da cellule morfologicamente differenziate, strettamente unite tra loro da \textbf{giunzioni specializzate} con scarsissima matrice extracellulare interposta. Le principali giunzioni cellulari sono:

\begin{itemize}
  \item \textbf{Desmosomi}: giunzioni di ancoraggio meccanico tra cellule adiacenti.
  \item \textbf{Giunzioni serrate} (\textit{tight junction}): sigillano lo spazio intercellulare, impedendo il passaggio di molecole tra le cellule.
  \item \textbf{Giunzioni comunicanti} (\textit{gap junction}): consentono la comunicazione diretta tra cellule adiacenti.
\end{itemize}

Il tessuto epiteliale assolve diverse funzioni: \textbf{protezione}, \textbf{scambio} e \textbf{lubrificazione} di superfici interne ed esterne. Queste funzioni vengono svolte mediante cellule altamente specializzate nelle attività di assorbimento, secrezione, trasporto, scambio di gas, scivolamento (tra due superfici) e ricezione sensoriale. Di conseguenza, la morfologia del tessuto di rivestimento si differenzia a seconda della funzione svolta.

I tessuti epiteliali si classificano in due grandi categorie:
\begin{itemize}
  \item \textbf{Epiteli di rivestimento}: ricoprono le superfici esterne ed interne del corpo.
  \item \textbf{Epiteli ghiandolari}: specializzati nella produzione e nel rilascio di secreti.
\end{itemize}

% ---------------------------------------------------------------------------
\section{La membrana basale}

Alla base di tutti gli epiteli è presente una \textbf{membrana basale} formata da proteoglicani, che separa l'epitelio dal tessuto connettivo sottostante. Al di sotto della membrana basale si trova tessuto connettivale reticolare lasso, ricco di \textbf{collagene di tipo III}.

La membrana basale è composta da due strati distinti:

\begin{itemize}
  \item \textbf{Lamina lucida}: strato più esterno, situato in prossimità della membrana plasmatica delle cellule epiteliali.
  \item \textbf{Lamina densa}: strato più profondo, secreto dalle cellule del connettivo sottostante; contiene collagene di tipo IV (nelle placche di ancoraggio) e collagene di tipo VII (nelle fibrille di ancoraggio).
\end{itemize}

La membrana basale è inoltre ricca di \textbf{molecole di adesione cellulare} (MAC) e \textbf{proteoglicani} (cemento intercellulare), che assicurano l'ancoraggio delle cellule epiteliali.

% ---------------------------------------------------------------------------
\section{Polarità del tessuto epiteliale}

Ad ogni epitelio corrisponde una \textbf{superficie esposta} (rivolta verso l'esterno del corpo o verso le cavità interne) e una \textbf{superficie basale} (dove l'epitelio prende contatto con le strutture circostanti). Le due superfici differiscono per struttura e funzione della membrana, poiché i componenti citoplasmatici cellulari non sono distribuiti uniformemente tra i due poli.

La superficie apicale può presentare strutture specializzate:
\begin{itemize}
  \item \textbf{Microvilli}: estroflessioni digitiformi che aumentano la superficie di assorbimento.
  \item \textbf{Ciglia}: appendici mobili per il trasporto di muco e particelle.
  \item \textbf{Stereociglia}: microvilli molto allungati, presenti nell'epididimo e nell'orecchio interno.
\end{itemize}

% ---------------------------------------------------------------------------
\section{Classificazione degli epiteli di rivestimento}

\subsection{In base al numero di strati cellulari}

\begin{itemize}
  \item \textbf{Epitelio monostratificato} (o semplice): un unico strato di cellule, tutte a contatto con la lamina basale.
  \item \textbf{Epitelio pseudostratificato}: le cellule poggiano tutte sulla lamina basale ma i nuclei si trovano a diverse altezze, dando l'apparenza di più strati.
  \item \textbf{Epitelio pluristratificato} (o composto): più strati sovrapposti; solo le cellule dello strato basale sono a contatto con la lamina basale.
\end{itemize}

\subsection{In base alla forma delle cellule}

\begin{itemize}
  \item \textbf{Pavimentoso}: cellule piatte, con larghezza e lunghezza nettamente superiori allo spessore.
  \item \textbf{Cubico}: cellule di forma approssimativamente cubica, con nucleo centrale sferico.
  \item \textbf{Cilindrico}: cellule più alte che larghe, con nucleo basale ovoidale.
\end{itemize}

\subsection{Epitelio cubico}

\textbf{Epitelio cubico semplice} --- Sede: ghiandole, dotti, tratti dei tubuli renali, ghiandola tiroidea. Funzioni: limitate funzioni protettive, secrezione e assorbimento.

\textbf{Epitelio cubico stratificato} --- Sede: ricopre alcuni dotti (raro). Funzioni: protezione, secrezione, assorbimento.

\subsection{Epitelio cilindrico}

\textbf{Epitelio cilindrico semplice} --- Sede: rivestimento dello stomaco, dell'intestino crasso, della colecisti, delle tube uterine e dei dotti collettori del rene. Funzioni: protezione, secrezione, assorbimento.

\textbf{Epitelio cilindrico stratificato} --- Sede: piccole zone di faringe, epiglottide, ano, dotti mammari, ghiandola salivare, uretra. Funzioni: protezione.

\subsection{Epitelio cilindrico ciliato pseudostratificato}

Sede: rivestimento della cavità nasale, della trachea e dei bronchi; porzioni dell'apparato riproduttivo maschile. Funzioni: protezione e secrezione. Le ciglia muovono il muco verso l'esterno, contribuendo all'eliminazione di particelle e agenti patogeni.

\subsection{Epitelio di transizione (urotelio)}

Sede: vescica urinaria, pelvi renale, ureteri. Funzioni: permette l'espansione e il ritorno dopo la distensione. Le cellule cambiano forma in base allo stato di riempimento dell'organo (cellule più piatte a vescica piena, più arrotondate a vescica vuota).

\subsection{Epitelio pavimentoso pluristratificato cheratinizzato --- l'epidermide}

Nell'\textbf{epidermide}, le cellule degli strati più superficiali perdono progressivamente i loro nuclei e il citoplasma viene occupato da una grande quantità di \textbf{cheratina}, una proteina fibrosa resistente. Le cellule superficiali diventano quindi non vitali e formano \textbf{squame cornee}, che costituiscono un'efficace barriera fisica contro traumi, disidratazione e agenti patogeni.

% ---------------------------------------------------------------------------
\section{Gli epiteli ghiandolari}

\subsection{Ghiandole esocrine ed endocrine}

Gli \textbf{epiteli ghiandolari} sono specializzati nella produzione di secreti. In base alla destinazione del secreto si distinguono:

\begin{itemize}
  \item \textbf{Epiteli ghiandolari esocrini}: cellule specializzate nella secrezione di sostanze che, attraverso un \textbf{dotto escretore}, vengono riversate sulla superficie esterna del corpo o all'interno di cavità. La porzione secernente, detta \textbf{adenomero}, rimane collegata all'epitelio tramite il dotto.
  \item \textbf{Epiteli ghiandolari endocrini}: cellule specializzate nella secrezione di sostanze (\textbf{ormoni}) che vengono riversate direttamente nei vasi sanguiferi. La porzione secernente si separa dalla superficie epiteliale e viene diffusamente vascolarizzata.
\end{itemize}

\subsection{Classificazione delle ghiandole esocrine pluricellulari}

Le ghiandole esocrine pluricellulari si classificano secondo quattro criteri:

\paragraph{A) In base alla posizione:}
\begin{itemize}
  \item \textbf{Intraparietali}: localizzate nella parete dell'organo che rivestono.
    \begin{itemize}
      \item Intraepiteliali: all'interno dello strato epiteliale.
      \item Esoepiteliali: al di sotto dell'epitelio (coriali o sottomucose).
    \end{itemize}
  \item \textbf{Extraparietali}: localizzate a distanza dall'organo che riforniscono (es.\ pancreas, fegato).
\end{itemize}

\paragraph{B) In base alla struttura del dotto e della porzione secernente:}

\noindent\textit{Ghiandole semplici} (dotto non ramificato):
\begin{itemize}
  \item \textbf{Tubulare semplice}: es.\ ghiandole intestinali (cripte di Lieberkühn).
  \item \textbf{Tubulo spirale semplice}: es.\ ghiandole sudoripare merocrine.
  \item \textbf{Tubulo ramificata semplice}: es.\ ghiandole gastriche, ghiandole mucose dell'esofago, della lingua e del duodeno.
  \item \textbf{Alveolare (acinosa) semplice}: non riscontrate nell'adulto; rappresentano uno stadio dello sviluppo delle ghiandole ramificate semplici.
  \item \textbf{Alveolare ramificata semplice}: es.\ ghiandole sebacee.
\end{itemize}

\noindent\textit{Ghiandole composte} (dotto ramificato):
\begin{itemize}
  \item \textbf{Tubulare composta}: es.\ ghiandole mucose della bocca, ghiandole bulbouretrali, tubuli seminiferi del testicolo.
  \item \textbf{Alveolare (acinosa) composta}: es.\ ghiandola mammaria.
  \item \textbf{Tubulo-alveolare composta}: es.\ ghiandole salivari, ghiandole del tratto respiratorio, pancreas.
\end{itemize}

\paragraph{C) In base al meccanismo di secrezione:}
\begin{itemize}
  \item \textbf{Merocrine}: il secreto viene rilasciato per \textbf{esocitosi} delle vescicole secretorie, senza perdita di citoplasma. È il meccanismo più comune (es.\ ghiandole salivari).
  \item \textbf{Apocrine}: il secreto si accumula nella porzione apicale della cellula, che viene poi distaccata e rilasciata; la cellula viene successivamente ripristinata (4 stadi: secrezione → decomposizione apicale → ripristino → nuovo ciclo). Es.\ ghiandola mammaria.
  \item \textbf{Olocrine}: le cellule accumulano il secreto fino a ``scoppiare'', rilasciando l'intero contenuto citoplasmatico. Le cellule perdute vengono rimpiazzate da mitosi delle cellule staminali basali (es.\ ghiandole sebacee).
\end{itemize}

\paragraph{D) In base al tipo di secreto:}
\begin{itemize}
  \item \textbf{Ghiandole sierose}: producono una secrezione acquosa, ricca di enzimi (es.\ amilasi salivare).
  \item \textbf{Ghiandole mucose}: secernono \textbf{mucine} (glicoproteine) che assorbono acqua formando muco viscoso (es.\ saliva mucosa).
  \item \textbf{Ghiandole miste}: contengono diversi tipi di cellule ghiandolari e producono sia secrezione sierosa che mucosa.
\end{itemize}

% ---------------------------------------------------------------------------
\section{Principali funzioni del tessuto epiteliale}

In sintesi, il tessuto epiteliale svolge quattro funzioni fondamentali:

\begin{itemize}
  \item \textbf{Protezione fisica}: barriera meccanica contro traumi, microrganismi e disidratazione.
  \item \textbf{Regolazione della permeabilità}: controllo selettivo del passaggio di sostanze tra i compartimenti.
  \item \textbf{Sensibilità}: ricezione di stimoli provenienti dall'ambiente esterno e interno.
  \item \textbf{Produzione di secreti specializzati}: sintesi e rilascio di enzimi, ormoni, muco e altre molecole biologicamente attive.
\end{itemize}
